\documentclass[11pt]{article}
\usepackage[T1]{fontenc}
\usepackage[utf8]{inputenc}
\usepackage{lmodern}
\usepackage[margin=1in]{geometry}
\usepackage{amsmath,amssymb,amsthm,mathtools}
\usepackage{booktabs,array,longtable}
\usepackage{enumitem}
\usepackage{hyperref}
\usepackage{xurl}

\IfFileExists{paper/release_info.tex}{
  \newcommand{\OPHPaperReleaseID}{r2011}
\newcommand{\OPHPaperReleaseDate}{August 1, 2026}
\newcommand{\OPHPaperReleaseBanner}{%
\begin{center}
\small\textbf{Paper release:} \texttt{\OPHPaperReleaseID}\quad\textbf{Released:} \OPHPaperReleaseDate
\end{center}
}
\newcommand{\PaperReleaseID}{\OPHPaperReleaseID}
\newcommand{\PaperReleaseDate}{\OPHPaperReleaseDate}
\newcommand{\PaperReleaseBanner}{%
\begin{center}
\small\textbf{Paper release:} \texttt{\PaperReleaseID}\quad\textbf{Released:} \PaperReleaseDate
\end{center}
}

}{
  \IfFileExists{../paper/release_info.tex}{
    \newcommand{\OPHPaperReleaseID}{r2011}
\newcommand{\OPHPaperReleaseDate}{August 1, 2026}
\newcommand{\OPHPaperReleaseBanner}{%
\begin{center}
\small\textbf{Paper release:} \texttt{\OPHPaperReleaseID}\quad\textbf{Released:} \OPHPaperReleaseDate
\end{center}
}
\newcommand{\PaperReleaseID}{\OPHPaperReleaseID}
\newcommand{\PaperReleaseDate}{\OPHPaperReleaseDate}
\newcommand{\PaperReleaseBanner}{%
\begin{center}
\small\textbf{Paper release:} \texttt{\PaperReleaseID}\quad\textbf{Released:} \PaperReleaseDate
\end{center}
}

  }{
    \newcommand{\OPHPaperReleaseID}{r1518}
    \newcommand{\OPHPaperReleaseDate}{July 7, 2026}
    \newcommand{\OPHPaperReleaseBanner}{%
      \begin{center}
      \small\textbf{Paper release:} \texttt{\OPHPaperReleaseID}\quad\textbf{Released:} \OPHPaperReleaseDate
      \end{center}}
  }
}

\hypersetup{colorlinks=true,linkcolor=blue,citecolor=blue,urlcolor=blue}
\setlength{\emergencystretch}{3em}
\setlength{\tabcolsep}{4pt}
\renewcommand{\arraystretch}{1.08}

\newtheorem{theorem}{Theorem}[section]
\newtheorem{definition}[theorem]{Definition}
\newtheorem{assumption}[theorem]{Assumption}
\newtheorem{target}[theorem]{Target}
\theoremstyle{remark}
\newtheorem{remark}[theorem]{Remark}
\theoremstyle{plain}

\newcommand{\pathref}[1]{\nolinkurl{#1}}
\newcommand{\SimSourceRepo}{\url{https://github.com/muellerberndt/oph-physics-sim}}
\newcommand{\SimEvidenceRepo}{\url{https://github.com/FloatingPragma/observer-patch-holography/tree/main/evidence}}

\title{Observer-Patch Holography Cosmological Vacuum and Structure Formation:\\
Finite Transfer Conditions, Fluctuation Ensembles, and Proto-Objects}
\author{B.\ M\"uller}
\date{\OPHPaperReleaseDate}

\begin{document}
\maketitle
\OPHPaperReleaseBanner

\begin{center}
\small\textbf{Author affiliation:} Bernhard Mueller, Pragma Research Inc.
\end{center}

\begin{abstract}
Observer-Patch Holography (OPH) cosmology needs a sharp distinction between a field that looks noisy and a
field that deserves to be called a physical vacuum. This paper sets that
boundary. A settled simulator state, a repair relaxation, or a plausible visual
fluctuation is only a diagnostic unless it comes from OPH source data with a
proper transfer construction and refinement check.

The same boundary organizes structure formation. Vacuum-like fields, primordial
seeds, and observable sky predictions are different claims. The paper keeps
those layers separate so a useful simulation cannot be mistaken for a completed
cosmological prediction.
\end{abstract}

\section{Scope}

This technical companion states the
vacuum and fluctuation boundary for finite-source cosmology, screen
microphysics, the compact results, and finite-patch simulations.
The simulator source is \SimSourceRepo, and its canonical public data
archives are indexed at \SimEvidenceRepo. The large finite repair archive is
a diagnostic whose physical-CMB and production-gravity gates are false; it
does not supply the source and transfer receipts required below.

\section{Finite Quotient Ensemble Theorem}

\input{../paper/tex_fragments/FINITE_QUOTIENT_ENSEMBLE_THEOREMS.tex}

\section{Vacuum Claim Boundary}

A simulated field is classified as seed noise, repair jitter, or a conventional
reference ensemble unless it satisfies the finite quotient ensemble
conditions. An OPH-native vacuum requires source data of the form
\[
\mathfrak S_r^E=(Q_r,m_r^0,J_r,V_r,a_{t,r})
\]
with \(J_r\), \(V_r\), and \(a_{t,r}\) derived from OPH source structure rather than from a target
law or a desired visual spectrum.  Accepted repair moves are generally directed and cannot simply
be reused as reversible Euclidean conductances.

\section{Structure-Formation Requirements}

\begin{longtable}{@{}p{0.28\textwidth}p{0.62\textwidth}@{}}
\toprule
Object & Required source-side label\\
\midrule
\endhead
\bottomrule
\endfoot
Fluctuating screen field & Seed noise or a conventional reference ensemble unless the quotient
ensemble conditions hold.\\
OPH vacuum field & Source Euclidean slab, ground-state, Doob, reflection-positivity,
and transfer-refinement conditions.\\
Proto-object worldlines & Diagnostic until production matter/particle and observer-facing object
population conditions hold.\\
Growth seeds & Primordial bridge, physical mode basis,
radial-null, positivity, forward-projection, and source-dilation or tomography
conditions.\\
Large-scale structure comparison & Diagnostic; source, transfer, and official likelihood/data
contracts are not supplied.\\
\end{longtable}

\section{Numerical Acceptance Conditions}

The finite-patch calculation reports the following five conditions separately:
\begin{itemize}[leftmargin=2em]
\item a tracially pointed quotient;
\item a source-derived Euclidean slab;
\item transfer stability under refinement;
\item schedule invariance of the stationary law;
\item pathwise invariance under partition changes.
\end{itemize}
A failed condition is part of the scientific result and excludes the physical
vacuum classification for a reference baseline or rendering field.

\section{Unsupplied Premises}

Finite OPH source data do not supply the reversible conductances \(J_r\), the
potential \(V_r\), or the slab thickness \(a_{t,r}\). Reflection positivity,
or an equivalent reconstruction theorem on the selected branch, is not
supplied. The bridge from a source-derived vacuum to primordial growth seeds
without cosmic microwave background or large-scale structure data is not
supplied. A physical proto-object or worldline requires conditions distinct
from those used for production particle claims.

\section*{AI Assistance Disclosure}
This research project used research-grade commercial models, including
Anthropic's Fable and OpenAI's GPT-5.6-Sol, for research support, software
development, editing, and synthesis. The authors are responsible for the
paper's claims, methods, and final text.

\end{document}
